% report_example.tex — structural reference for the /synthesize report template.
% This is a compile/structure demo, not a live-verified synthesis output: a real
% /synthesize run writes reports/<topic_slug>/report.tex from sources actually
% fetched during that run, per 04-citation-rules.md. See 03-report-templates.md
% for the section-by-section rules this file demonstrates.
\documentclass[11pt]{article}

\usepackage[margin=1in]{geometry}
\usepackage{hyperref}
% natbib's citation mode must match \bibliographystyle: [numbers] for numbered
% styles (ieeetr, plain); no option (default author-year) for apalike/plainnat.
% See 04-citation-rules.md's style table.
\usepackage[numbers,sort&compress]{natbib}
\usepackage{booktabs}

\hypersetup{
  colorlinks=true,
  linkcolor=black,
  citecolor=blue,
  urlcolor=blue,
}

\title{Retrieval-Augmented Generation: A Synthesis}
\author{Prepared by AI Research Assistant}
\date{\today}

\begin{document}
\maketitle

\noindent\textit{Search scope: arXiv (cs.CL, cs.LG) and Semantic Scholar, queried for
``retrieval-augmented generation'' and related terms. This example demonstrates report
structure and LaTeX/BibTeX mechanics; it is not a live-verified output of a
\texttt{/synthesize} run.}

\begin{abstract}
Retrieval-augmented generation (RAG) combines a parametric language model with a
non-parametric retrieval mechanism over an external corpus, letting a generation model
condition on retrieved evidence rather than relying solely on knowledge encoded in its
weights. Early approaches established the core recipe: encode a query, retrieve
relevant passages with a dense retriever, and condition generation on the retrieved
context~\citep{lewis2020rag,karpukhin2020dpr,guu2020realm}. Subsequent work explored
how retrieved passages are fused into generation~\citep{izacard2021fid}. The main open
questions concern how retrieval and generation should be jointly trained, and how to
evaluate factual grounding rather than just downstream task accuracy.
\end{abstract}

\section{Background}
A parametric language model stores factual knowledge implicitly in its weights, which
makes updating that knowledge or attributing a specific claim to a source difficult.
Retrieval-augmented approaches instead give the model access to an external,
inspectable corpus at inference time, addressing both problems: the corpus can be
updated without retraining, and retrieved passages provide provenance for generated
claims.

\section{Approaches}

\subsection{Dense Retrieval Foundations}
Dense Passage Retrieval~\citep{karpukhin2020dpr} showed that a dual-encoder retriever,
trained with in-batch negatives, substantially outperforms sparse lexical retrieval
(e.g.\ BM25) for open-domain question answering. This dense-retrieval recipe became the
standard retrieval component for the generation-focused work that followed.

\subsection{End-to-End Differentiable Retrieval}
REALM~\citep{guu2020realm} and the original RAG
formulation~\citep{lewis2020rag} both make the retrieval step differentiable with
respect to the end task, so the retriever can be fine-tuned jointly with the generator
rather than treated as a frozen upstream component.
RAG in particular formulates two variants: RAG-Sequence, which conditions on the same
retrieved passages for the whole output, and RAG-Token, which can attend to different
passages at each generated token.

\subsection{Fusion Strategies for Generation}
Fusion-in-Decoder~\citep{izacard2021fid} processes each retrieved passage
independently in the encoder, then fuses all of them in a single decoder pass. This
scales more gracefully to a larger number of retrieved passages than RAG's approach of
marginalizing over passages, at the cost of encoder compute proportional to the number
of passages retrieved.

\section{Technical Findings (Plain Language)}
\textbf{Dense retrieval, in plain terms:} instead of matching a search query against
documents by shared keywords (the old way), dense retrieval turns both the query and
every candidate document into a list of numbers (a vector) using a neural network, and
finds documents whose numbers are ``close'' to the query's numbers. This is why it
beats keyword matching: it can find a relevant passage even when it doesn't share any
exact words with the query, because it's matching on meaning, not vocabulary.

\textbf{RAG and REALM, in plain terms:} a normal language model only knows what it
memorized during training. RAG and REALM bolt on a second step - before answering,
first fetch a handful of relevant text passages from an external database, then hand
those passages \emph{along with} the question to the language model. The model's job
becomes ``answer using these specific passages,'' not ``recall everything from
training.'' REALM builds this fetch-then-answer habit in from the start of training;
RAG adds it on top of an already-trained model. The practical upshot: you can update
what the model ``knows'' by updating the external database, no retraining required.

\textbf{Fusion-in-Decoder, in plain terms:} RAG has to average across every fetched
passage in a single pass, which gets harder to do well as you fetch more passages.
Fusion-in-Decoder instead reads each passage separately first, then combines
everything only at the very last step of generating the answer. This lets it use more
passages without the averaging problem getting worse - the tradeoff is that reading
more passages separately costs more compute.

\section{Comparison of Approaches}

\begin{table}[h]
\centering
\begin{tabular}{@{}p{2.6cm}p{3.4cm}p{4cm}p{4cm}@{}}
\toprule
\textbf{Approach} & \textbf{Key Paper} & \textbf{Strengths} & \textbf{Limitations} \\
\midrule
RAG & Lewis et al.~\citep{lewis2020rag} & Joint retriever/generator fine-tuning; two token/sequence variants & Retrieval quality bounded by a single dense retriever pass \\
REALM & Guu et al.~\citep{guu2020realm} & Retrieval pre-trained jointly with masked LM objective & Pre-training cost; less flexible at fine-tuning time \\
FiD & Izacard \& Grave~\citep{izacard2021fid} & Scales well to many retrieved passages via per-passage encoding & Encoder cost grows linearly with passages retrieved \\
\bottomrule
\end{tabular}
\caption{Representative approaches to retrieval-augmented generation.}
\end{table}

\section{Open Questions}
\begin{itemize}
  \item \textbf{Joint vs.\ frozen retrieval:} it remains unsettled whether jointly
    fine-tuning the retriever consistently outperforms a strong frozen retriever
    (e.g.\ DPR) paired with a well-tuned generator, or whether gains are
    task-dependent.
  \item \textbf{Evaluating faithfulness:} downstream task accuracy does not directly
    measure whether generated text is actually grounded in the retrieved passages;
    the surveyed works report task metrics rather than faithfulness metrics.
  \item \textbf{Passage count vs.\ compute:} FiD's linear encoder cost in the number of
    passages trades off against RAG's fixed-cost marginalization; no work surveyed
    here directly compares the two under matched compute budgets.
\end{itemize}

\section{Evidence Basis}
Per source, how its content was obtained and what a reader should weigh when relying
on it. Stored per entry in \texttt{references.bib}'s \texttt{evidencebasis} field
(a custom field no stock \texttt{.bst} style prints - see \texttt{04-citation-rules.md})
and collected here as a table rather than left to render inline after each reference.

\begin{table}[h]
\centering
\begin{tabular}{@{}p{2.0cm}p{5.1cm}p{2.2cm}p{3.2cm}@{}}
\toprule
\textbf{Source} & \textbf{Evidence Basis} & \textbf{Disclosure} & \textbf{Caveat} \\
\midrule
Lewis et al.~\citep{lewis2020rag} & Illustrative entry (structural demo, not a live fetch) & academic & None \\
Karpukhin et al.~\citep{karpukhin2020dpr} & Illustrative entry (structural demo, not a live fetch) & academic & None \\
Guu et al.~\citep{guu2020realm} & Illustrative entry (structural demo, not a live fetch) & academic & None \\
Izacard \& Grave~\citep{izacard2021fid} & Illustrative entry (structural demo, not a live fetch) & academic & None \\
\bottomrule
\end{tabular}
\caption{Evidence basis per source - a real \texttt{/synthesize} run states an actual
full-text/abstract-only basis and route per source, per \texttt{03-report-templates.md}.}
\end{table}

\bibliographystyle{ieeetr}
\bibliography{references}

\end{document}
